% ============================================================
% pure 报告 — 规范模板（xelatex + ctex）
% 主题: DDalphaAMG 核心部分穷尽剖析
% 编译: xelatex 两遍; 交付前 Overfull / Float too large 均为 0
% ============================================================
\documentclass[UTF8]{ctexart}
\usepackage[margin=2.5cm]{geometry}
\usepackage{graphicx}
\usepackage{amsmath}
\usepackage{amssymb}
\usepackage{microtype}
\usepackage{listings}
\usepackage{xcolor}
\usepackage{booktabs}
\usepackage{longtable}
\usepackage{xltabular}
\usepackage{tabularx}
\usepackage{hyperref}
\usepackage{fancyhdr}
\usepackage[most]{tcolorbox}
\usepackage{titlesec}

\definecolor{accent}{HTML}{1F4E79}
\definecolor{evidence}{HTML}{1E8449}
\definecolor{reference}{HTML}{8C6D1F}
\definecolor{conclusion}{HTML}{8B1A1A}
\definecolor{codebg}{HTML}{F4F6F7}
\definecolor{algo}{HTML}{E67E22}
\definecolor{phys}{HTML}{6C3483}
\definecolor{map}{HTML}{C2185B}
\definecolor{warn}{HTML}{D35400}
\definecolor{hl}{HTML}{FDF6E3}

\setlength{\emergencystretch}{3em}
\hypersetup{colorlinks=true,linkcolor=accent,urlcolor=map}

\titleformat{\section}{\Large\bfseries\color{accent}}{\thesection}{1em}{}
\titleformat{\subsection}{\large\bfseries\color{phys}}{\thesubsection}{1em}{}
\titleformat{\subsubsection}{\normalsize\bfseries\color{phys!70!black}}{\thesubsubsection}{1em}{}

\newtcolorbox{conclbox}{breakable,colback=conclusion!6,colframe=conclusion,title=结论}
\newtcolorbox{refbox}{breakable,colback=reference!6,colframe=reference,title=参考背景}
\newtcolorbox{warnbox}{breakable,colback=warn!6,colframe=warn,title=局限与风险}
\newtcolorbox{keybox}{breakable,colback=hl,colframe=accent,title=要点}

\lstset{
  basicstyle=\ttfamily\footnotesize,
  backgroundcolor=\color{codebg},
  frame=single,
  language=C,
  firstnumber=auto,
  keywordstyle=\color{accent}\bfseries,
  commentstyle=\color{evidence!70!black},
  stringstyle=\color{phys},
  breaklines=true,
  breakatwhitespace=false,
  postbreak=\mbox{\textcolor{accent}{$\hookrightarrow$}\space},
  keepspaces=true,
  columns=flexible,
  numbers=left,numberstyle=\tiny\color{phys!60!black},
}

\title{DDalphaAMG 核心部分穷尽剖析报告}
\author{opencode pure}
\date{2026-08-17}

\begin{document}
\maketitle

\begin{abstract}
本报告对 \url{/root/PyQCU/refer/git-rep/DDalphaAMG} 的核心部分进行穷尽剖析。
项目性质判定为\textbf{代码-物理交叉向}：聚合型代数多重网格（AMG）求解 Wilson--Clover 狄拉克
方程。穷尽枚举 12 个核心候选（C1--C12），其中\textbf{深挖 7 / 浅析 3 / 排除 2}：
深挖部分含物理图像、公式推导链、代码-物理映射表与证据（\texttt{文件:行号}）。
独特竞争力定位：聚合 AMG + SAP 平滑 + 自适应 setup 的格子无关收敛率、sed 精度实例化与
CUDA/SSE 多后端。注意：本快照缺 CPU \texttt{.c} 实现与 Makefile，证据以
\texttt{*\_generic.h} 声明与 \texttt{src/gpu/*.cu} 为准；非独立 git 仓库。
\end{abstract}

\tableofcontents

\section{项目定位与核心部分清单}

\subsection{项目性质判定}
\begin{keybox}
\textbf{性质判定：交叉向。}本库是格点 QCD 求解器：算法代码（AMG/SAP/Krylov）为核心载体，
其物理对象明确（Wilson--Clover 算子、规范 link、clover 项、test vectors），存在密集的
代码符号 $\leftrightarrow$ 物理量对应关系（见 C1--C12）。判定依据：\texttt{src/} 129 文件全为算法头与 CUDA
内核，\texttt{doc/user\_doc.tex:21-26} 明确求解 $D_W(U,m_0)\psi=\eta$。
\end{keybox}

\subsection{核心部分总清单（穷尽枚举）}
\begin{xltabular}{\textwidth}{lXXX}
\toprule
\textcolor{map}{编号} & 候选 & 三态 & 理由 \\
\midrule
\endhead
C1 & Wilson--Clover 算子 & 深挖 & 求解方程核心，代码-物理映射密集 \\
C2 & 偶奇（Schur 补）预处理 & 深挖 & 收敛加速核心机制 \\
C3 & 自适应 setup & 深挖 & 格子无关收敛的关键 \\
C4 & 几何聚合与插值算子 & 深挖 & AMG 构造核心 \\
C5 & 粗算子三重积 $P^\dagger D P$ & 深挖 & 粗格正确性核心 \\
C6 & SAP 平滑器 & 深挖 & 平滑器核心 \\
C7 & V/K-cycle 与 GCR-ODR & 深挖 & 循环策略核心 \\
C8 & FGMRES 与混合精度 & 浅析 & 外层求解器，支撑性 \\
C9 & ghost 通信重叠 & 浅析 & 性能优化，非算法核心 \\
C10 & SIMD/CUDA 后端 & 浅析 & 性能移植，架构独特但非物理核心 \\
C11 & sed 精度实例化 & 排除 & 工程机制，已在 analy 覆盖 \\
C12 & 配置/INI 解析 & 排除 & 通用工程，无独特性 \\
\bottomrule
\caption{核心部分清单三态（数据来源: 全库索引实测）}
\end{xltabular}

\section{核心部分深度剖析}

\subsection{C1 Wilson--Clover 算子}
\textbf{物理图像：}Wilson 费米子作用量在格点上的 Dirac 算子 $D_W$ 由"跳跃项 + 质量项 +
\begin{xltabular}{\textwidth}{lXX}
\toprule
\textcolor{map}{代码符号} & 物理/算法对象 & 含义 \\
\midrule
\endhead
\texttt{D（36i+9mu）} & 规范 link $U_\mu(x)$ & 每方向 9 复分量跳跃矩阵 \\
\texttt{op->clover（42i）} & Clover 项 & 12$\times$12 块（对角+上三角，含 $4+m_0$）\\
\texttt{prp\_mu/prn\_mu} & spin 投影 & 投影 $\frac12(1\mp\gamma_\mu)$ \\
\texttt{gamma5\_PRECISION} & $\gamma_5$-Hermite & $D^\dagger=\gamma_5 D\gamma_5$\\
\texttt{csw / m0} & $c_{sw}$, $m_0$ & clover 系数、裸质量 \\
\bottomrule
\caption{映射表 C1（数据来源: src/*.h 与 doc/user\_doc.tex 实测）}
\end{xltabular}
\textbf{推导链：}Wilson 算子
\begin{equation}\label{eq:wilson}
D_W(U,m_0) = (4+m_0)\mathbb{1} - \frac12 \sum_\mu\left[(1-\gamma_\mu)U_\mu(x)\delta_{x+\hat\mu,y} + (1+\gamma_\mu)U_\mu^\dagger(x-\hat\mu)\delta_{x-\hat\mu,y}\right].
\end{equation}
Clover 修正加 $c_{sw}\sigma_{\mu\nu}F_{\mu\nu}(x)$ 项（\nolinkurl{src/dirac_generic.h:636-741}
site clover 12$\times$12 块作用）。GPU 实现见 \nolinkurl{src/gpu/cuda_dirac_generic.cu:171-211}
（\nolinkurl{cuda_block_d_plus_clover_PRECISION}）。
\textbf{极限校验：}$m_0=-4$（$\kappa\to\infty$）时对角项消失，对应无质量极限；$\gamma_5$-Hermite
对称性保证粗算子块结构 $A=A^\dagger$、$D=D^\dagger$、$C=-B^\dagger$
（\nolinkurl{src/coarse_operator_generic.h:124-127}）。

\subsection{C2 偶奇（Schur 补）预处理}
\textbf{物理图像：}格点按奇偶（checkerboard）着色后，Dirac 算子呈块结构，Schur 补
$\hat D = D_{oo} - D_{oe}D_{ee}^{-1}D_{eo}$ 条件数远小于 $D$，收敛加速约 $\sqrt{2}$ 倍。
\textbf{代码定位：}\nolinkurl{src/oddeven_generic.h:27-64}（\texttt{hopping\_term\_PRECISION}、
\nolinkurl{apply_schur_complement_PRECISION_cpu}、\texttt{block\_*} 块级 Schur）、
\nolinkurl{src/coarse_oddeven_generic.h:27-53}（粗层偶奇）。
\textbf{推导链：}令 $D=\begin{pmatrix} D_{ee} & D_{eo}\\ D_{oe} & D_{oo}\end{pmatrix}$，
Schur 补作用为 $\hat D\psi_o = \eta_o - D_{oe}D_{ee}^{-1}\eta_e$。
\textbf{极限校验：}自由场 $U=\mathbb{1}$ 时 Schur 补保持 Wilson 谱结构；GPU 最大实现
\nolinkurl{src/gpu/cuda_oddeven_generic.cu}（4532 行）为块级偶奇求解主体。

\subsection{C3 自适应 setup}
\textbf{物理图像：}以随机初值生成 test vectors 并平滑，捕捉近零模子空间，构造插值算子使
粗格继承低频信息。\textbf{代码定位：}\nolinkurl{src/setup_generic.h:27-32}
（\nolinkurl{iterative_PRECISION_setup}、\nolinkurl{interpolation_PRECISION_define}、
\nolinkurl{inv_iter_inv_fcycle_PRECISION}）。
\textbf{推导链：}每聚合 $n$ 个 test vector $\rightarrow$ 聚合内 Gram--Schmidt 正交化
（\nolinkurl{src/linalg_generic.h:102-143}）$\rightarrow$ 插值 $P$。
\textbf{为什么自适应：}静态插值对复杂规范场失效；自适应迭代（平滑-正交化-再插值）使
$P$ 逼近近零模空间，实现格子无关收敛。

\subsection{C4 几何聚合与插值算子}
\textbf{物理图像：}粗格点 = 细格点的聚合（aggregate），聚合尺寸 $agg_i(\mu)$ 控制粗化比。
\textbf{代码定位：}\nolinkurl{src/coarsening_generic.h:25-29}
（\nolinkurl{coarsening_index_table_PRECISION_define}）、
\nolinkurl{src/interpolation_generic.h:32-36}（\texttt{interpolate}/\texttt{restrict}=$P^\dagger$）。
\textbf{推导链：}细格 test vectors 按聚合分组 $\rightarrow$ $P$ 由聚合内 test vectors 构成，
限制为 $P^\dagger$。几何约束见 \texttt{doc/user\_doc.tex:53-65}（$agg_0=4$、$agg_i=2$）。

\subsection{C5 粗算子三重积 $P^\dagger D P$}
\textbf{物理图像：}Galerkin 粗化：粗算子 $D_c=P^\dagger D P$ 保持低频谱，粗格解经 $P$ 延拓
回细格。\textbf{代码定位：}\nolinkurl{src/coarse_operator_generic.h:31}
（\nolinkurl{coarse_operator_PRECISION_setup} 三重积）、\texttt{:45}
（\texttt{apply\_coarse\_operator}）、\texttt{:119-336}（\texttt{coarse\_hopp} 四块列存
A/B/C/D 乘法）。\textbf{推导链：}
\begin{equation}\label{eq:galerkin}
D_c = P^\dagger D P,\quad D_c^{\dagger}=\gamma_5 D_c\gamma_5 .
\end{equation}
\textbf{极限校验：}$D$ 自伴（$\gamma_5$-Hermite）保证 $D_c$ 块结构 $A=A^\dagger$、$D=D^\dagger$、
$C=-B^\dagger$（\nolinkurl{src/coarse_operator_generic.h:124-127}），粗层可用更小 Krylov 子空间。

\subsection{C6 SAP 平滑器}
\textbf{物理图像：}Schwarz 交替过程（SAP）将格点分为块，块内精确求解、块间交替更新，
作为 AMG 预/后平滑器消除高频误差。\textbf{代码定位：}\nolinkurl{src/schwarz_generic.h:49-56}
（\texttt{additive\_schwarz}、\texttt{red\_black\_schwarz}、
\texttt{sixteen\_color\_schwarz}）、\texttt{:36-61}（\texttt{schwarz\_setup}、
\texttt{connect\_link}）。
\textbf{推导链：}块分解 $\Omega=\cup_k\Omega_k$，加法 Schwarz 迭代
$\phi \leftarrow \phi + \sum_k R_k^T A_k^{-1} R_k(b - A\phi)$。
\textbf{为什么多着色：}红黑（2 色）保证块间无依赖可并行；16 色用于更细粒度并行。

\subsection{C7 V/K-cycle 与 GCR-ODR}
\textbf{物理图像：}V-cycle：细格预平滑 $\rightarrow$ 限制 $\rightarrow$ 递归粗解 $\rightarrow$
延拓 $\rightarrow$ 校正 $\rightarrow$ 后平滑；K-cycle：每层用 FGMRES 包裹两格方法
（\texttt{doc/user\_doc.tex:85}）。\textbf{代码定位：}\texttt{src/vcycle\_generic.h:38}
（\texttt{vcycle\_PRECISION}）、\nolinkurl{src/gcrodr_generic.h:25-32}
（\texttt{flgcrodr\_PRECISION} 最粗层 GCR-ODR 回收子空间）。
\textbf{为什么 K-cycle：}V-cycle 粗层收敛不足时，K-cycle 以 FGMRES 重新组合，提升收敛率
而保持成本可控。

\section{独特性与竞争力评估}
\begin{xltabular}{\textwidth}{lXX}
\toprule
\textcolor{algo}{核心} & 独特竞争力 & 对比朴素实现 \\
\midrule
\endhead
聚合 AMG + SAP & 格子无关收敛率（文献证实） & 纯 Krylov 收敛随体积恶化 \\
自适应 setup & 近零模子空间自动捕捉 & 静态插值对复杂场失效 \\
sed 精度实例化 & double/float 双实例共存 & 手工复制两份代码 \\
CUDA/SSE 多后端 & 单算法多架构 & 平台绑定实现 \\
$\gamma_5$-Hermite 块结构 & 粗算子压缩存储 & 全稠密粗算子 \\
\bottomrule
\caption{独特性评估（数据来源: 源码与文献）}
\end{xltabular}

\section{关键片段与公式}
\begin{lstlisting}[language=C]
// eta += D*phi, D stored columnwise
static inline void mv_PRECISION(const vector_PRECISION eta,
                                const complex_PRECISION *D,
                                const vector_PRECISION phi,
                                const register int n) {
  register int i, j, k=0;
  for (i=0; i<n; i++)
    for (j=0; j<n; j++, k++)
      eta[j] += D[k]*phi[i];
}
\end{lstlisting}
参考源：\nolinkurl{src/coarse_operator_generic.h:63-70}（列存粗算子乘）

核心公式汇总：
\begin{align}
D_W &= (4+m_0)\mathbb{1} - \tfrac12\sum_\mu[(1-\gamma_\mu)U_\mu\delta_{+}+(1+\gamma_\mu)U_\mu^\dagger\delta_{-}], \label{eq:w2}\\
D_c &= P^\dagger D P \quad(\text{Galerkin}),\quad D_c^{\dagger}=\gamma_5 D_c\gamma_5, \label{eq:g2}\\
m_0 &= 1/(2\kappa)-4,\quad np=\prod_\mu \texttt{d0 global}/\texttt{d0 local}. \label{eq:m0}
\end{align}

\section{参考源清单}
{\small
\begin{xltabular}{\textwidth}{XXX}
\toprule
\textcolor{evidence}{参考源} & 类型 & 内容/作用 \\
\midrule
\endhead
\texttt{\nolinkurl{src/dirac_generic.h:636-741}} & 仓库 & site clover 12$\times$12 块 \\
\texttt{\nolinkurl{src/gpu/cuda_dirac_generic.cu:171-211}} & 仓库 & GPU Wilson+Clover \\
\texttt{\nolinkurl{src/oddeven_generic.h:27-64}} & 仓库 & 偶奇/Schur 补 \\
\texttt{\nolinkurl{src/coarse_oddeven_generic.h:27-53}} & 仓库 & 粗层偶奇 \\
\texttt{\nolinkurl{src/setup_generic.h:27-32}} & 仓库 & 自适应 setup \\
\texttt{\nolinkurl{src/linalg_generic.h:102-143}} & 仓库 & Gram--Schmidt \\
\texttt{\nolinkurl{src/coarsening_generic.h:25-29}} & 仓库 & 聚合索引表 \\
\texttt{\nolinkurl{src/interpolation_generic.h:32-36}} & 仓库 & 插值/限制 \\
\texttt{\nolinkurl{src/coarse_operator_generic.h:31-37}} & 仓库 & 三重积 setup \\
\texttt{\nolinkurl{src/coarse_operator_generic.h:119-127}} & 仓库 & coarse\_hopp 块结构 \\
\texttt{\nolinkurl{src/schwarz_generic.h:49-56}} & 仓库 & SAP 三变体 \\
\texttt{\nolinkurl{src/vcycle_generic.h:38}} & 仓库 & V-cycle \\
\texttt{\nolinkurl{src/gcrodr_generic.h:25-32}} & 仓库 & GCR-ODR \\
\texttt{\nolinkurl{doc/user_doc.tex:21-26,53-85}} & 仓库 & 方程/几何/K-cycle \\
\texttt{\nolinkurl{doc/user_doc.bib:1-29}} & 仓库 & 文献（Frommer 2013/2014）\\
\texttt{\nolinkurl{src/coarse_operator_generic.h:63-70}} & 仓库 & mv 列存乘 \\
\bottomrule
\end{xltabular}
}

\section{结论}
\begin{conclbox}
DDalphaAMG 的核心竞争力是\textbf{聚合 AMG + SAP + 自适应 setup}的组合：以 $\gamma_5$-Hermite
块结构压缩粗算子、以 sed 精度实例化与 CUDA/SSE 分派换取多后端可移植，实现格子无关收敛率。
穷尽枚举 12 个候选，深挖 7 项均有物理图像与代码-物理映射证据。
\end{conclbox}
\begin{warnbox}
\textcolor{warn}{局限：}本快照缺 CPU \texttt{.c} 实现与 Makefile，部分推导以
\texttt{*\_generic.h} 声明与 GPU \texttt{.cu} 为证据；非独立 git 仓库；未实测运行。
\end{warnbox}

\end{document}